\subsection{EHZ capacity}\label{subsec:preliminaries-smooth-symplectic-geometry}

For smooth convex bodies, the Ekeland--Hofer and Hofer--Zehnder capacities
coincide. The common value is the least action needed to close up a
characteristic trajectory on the boundary. We use the resulting EHZ capacity in
the form
\begin{equation}\label{eq:preliminaries-ehz-capacity}
  c_{\mathrm{EHZ}}(K)
  = \min\{A(\gamma):\gamma
    \text{ is a closed characteristic on }\partial K\}.
\end{equation}
Equivalently, the minimum may be taken over Reeb orbits. The contact
normalization only chooses the parametrization of each characteristic, while
the action is a line integral and is invariant under orientation-preserving
reparametrization. For a contact-normalized representative, \(A(\gamma)\) is
also the period. The minimum exists and is positive; for the thesis we use this
standard existence result as background \cite{Rabinowitz1978}. For nonsmooth
convex bodies, and in particular for polytopes, the same capacity is obtained
by approximation and by minimizing over the generalized characteristics of
Definition~\ref{def:preliminaries-generalized-characteristic}. Their finite
facet description on a polytope is derived in
Section~\ref{sec:generalized-reeb-orbits-polytopes}.

As a symplectic capacity on convex bodies, \(c_{\mathrm{EHZ}}\) is monotone
under inclusion, invariant under symplectomorphisms, and conformal:
\[
  c_{\mathrm{EHZ}}(\lambda K)=\lambda^2c_{\mathrm{EHZ}}(K)
  \quad(\lambda>0).
\]
It is normalized by
\[
  c_{\mathrm{EHZ}}(B^4(r))=c_{\mathrm{EHZ}}(Z^4(r))=\pi r^2,
\]
where \(B^4(r)\) is the radius-\(r\) ball and
\(Z^4(r)=B^2(r)\times\mathbb{R}^2\) is the symplectic cylinder. These axioms
are used only as background; the finite computations below use the
minimum-action and quadratic-program descriptions.

Throughout the thesis, \(\operatorname{vol}_4\) denotes four-dimensional
Lebesgue measure in the coordinate order \((q_1,q_2,p_1,p_2)\). In dimension
four we measure the normalized capacity-to-volume ratio by
\begin{equation}\label{eq:preliminaries-systolic-ratio}
  \operatorname{sys}(K)
  \coloneqq
  \frac{c_{\mathrm{EHZ}}(K)^2}{2\operatorname{vol}_4(K)}.
\end{equation}
With this normalization, Viterbo's conjecture for four-dimensional convex
bodies says
\[
  \operatorname{sys}(K)\leq 1,
  \qquad\text{equivalently}\qquad
  c_{\mathrm{EHZ}}(K)^2\leq 2\operatorname{vol}_4(K)
\]
\cite{Viterbo2000}. Later sections use
\(\operatorname{sys}\) for this ratio for polytopes as well, via the same
capacity value discussed above.

\begin{definition}[Systolic-ratio symmetries]
\label{def:preliminaries-systolic-ratio-symmetries}
Let
\[
  G_{\operatorname{sys}}
  =
  \{x\mapsto b+\rho Mx:
    b\in\mathbb R^4,\ \rho>0,\ M\in\operatorname{Sp}(4,\mathbb R)\}.
\]
This is the group of affine transformations generated by translations,
positive scalings, and linear symplectic transformations.
\end{definition}

\begin{lemma}[Systolic-ratio invariance]
\label{lem:preliminaries-systolic-ratio-invariance}
For every convex body \(K\subset\mathbb R^4\) and every
\(g\in G_{\operatorname{sys}}\),
\[
  \operatorname{sys}(gK)=\operatorname{sys}(K).
\]
\end{lemma}

\begin{proof}
Translations and linear symplectic transformations are symplectomorphisms, so
they preserve \(c_{\mathrm{EHZ}}\). They also preserve four-dimensional
volume, because translations do not change volume and every symplectic matrix
has determinant \(1\). A positive scaling by \(\rho\) gives
\[
  c_{\mathrm{EHZ}}(\rho K)=\rho^2c_{\mathrm{EHZ}}(K),
  \qquad
  \operatorname{vol}_4(\rho K)=\rho^4\operatorname{vol}_4(K).
\]
Substituting these identities in
\eqref{eq:preliminaries-systolic-ratio} shows that the ratio is unchanged.
\end{proof}

This is a geometric symmetry group acting on unlabelled convex bodies. In a
labelled coordinate chart, the same unlabelled body may also be represented
after relabelling facets or inequalities; that relabelling ambiguity is a
coordinate issue, not an additional element of \(G_{\operatorname{sys}}\).

\begin{proposition}[Continuity]\label{prop:preliminaries-sys-hausdorff-continuity}
Let \(K_t\) be a family of four-dimensional convex bodies which is continuous
in the Hausdorff topology. Then \(c_{\mathrm{EHZ}}(K_t)\) is continuous in
\(t\). If, moreover, \(t\mapsto\operatorname{vol}_4(K_t)\) is positive,
finite, and continuous, then the systolic ratio of the family is continuous.
\end{proposition}

\begin{proof}
The EHZ capacity is \(C^0\)-continuous on convex bodies with respect to the
Hausdorff metric \cite{CH2021}. The assertion for \(\operatorname{sys}\) is
then ordinary continuity of the quotient in
\eqref{eq:preliminaries-systolic-ratio}. In the applications below the volume
hypothesis is usually automatic, since four-volume is continuous under
Hausdorff convergence of convex bodies; for the rotated-polygon family it is
even constant.
\end{proof}
